%%%%%%%%%%%%%%%%%%%%%%%%%%%%%%%%%%%%%%%%%%%%%%%%%%%%%%%%%%%%%%%%%%%%%%
% Lớp 9 | Chương VI. Hàm số y = ax² (a ≠ 0). Phương trình bậc hai một ẩn | Hàm số y = ax² (a ≠ 0)
% Dạng 9D6?18-2: Vẽ đồ thị, đọc đồ thị
%%%%%%%%%%%%%%%%%%%%%%%%%%%%%%%%%%%%%%%%%%%%%%%%%%%%%%%%%%%%%%%%%%%%%%


% nguồn: 2026-11-14_09_de_tuyen_sinh_10_tphcm_2025_2026 (nhập 14/11/2026)
\begin{bt}%[9D6H18-2]%[ma:f280a39f]
% nguon: Đề thi chính thức Tuyển sinh 10 TP.HCM năm học 2025-2026 – Bài 1
Cho hàm số $y = \dfrac{x^2}{2}$ có đồ thị $(P)$.
\begin{enumerate}[a)]
	\item Vẽ đồ thị $(P)$ trên hệ trục toạ độ.
	\item Tìm toạ độ các điểm thuộc $(P)$ có tung độ bằng $18$.
\end{enumerate}
\loigiai{
\begin{enumerate}[a)]
	\item Bảng giá trị tương ứng giữa $x$ và $y$ của hàm số $y = \dfrac{x^2}{2}$:
	\begin{center}
	\begin{tabular}{|c|c|c|c|c|c|}
	\hline
	$x$ & $-4$ & $-2$ & $0$ & $2$ & $4$ \\
	\hline
	$y = \dfrac{x^2}{2}$ & $8$ & $2$ & $0$ & $2$ & $8$ \\
	\hline
	\end{tabular}
	\end{center}
	Đồ thị $(P)$ là một parabol đi qua gốc toạ độ $O(0; 0)$, nhận trục tung $Oy$ làm trục đối xứng và nằm phía trên trục hoành.
	\item Điểm thuộc $(P)$ có tung độ bằng $18$ thoả mãn phương trình:
	$18 = \dfrac{x^2}{2} \Leftrightarrow x^2 = 36 \Leftrightarrow x = 6$ hoặc $x = -6$.\\
	Vậy toạ độ các điểm thoả mãn yêu cầu đề bài là $(6; 18)$ và $(-6; 18)$.
\end{enumerate}
}
\end{bt}

% nguồn: 2026-11-14_09_de_minh_hoa_tuyen_sinh_10_tphcm (nhập 14/11/2026)
\begin{bt}%[9D6H18-2]%[ma:f13c3ea3]
% nguon: Đề thi tham khảo Tuyển sinh 10 TP.HCM – Phụ lục Sở GD&ĐT TP.HCM – Bài 1
Cho hàm số $y = \dfrac{x^2}{2}$ có đồ thị $(P)$.
\begin{enumerate}[a)]
	\item Vẽ đồ thị $(P)$ của hàm số trên.
	\item Tìm những điểm $M$ thuộc $(P)$ có tung độ và hoành độ bằng nhau.
\end{enumerate}
\loigiai{
\begin{enumerate}[a)]
	\item Bảng giá trị tương ứng giữa $x$ và $y$ của hàm số $y = \dfrac{x^2}{2}$:
	\begin{center}
	\begin{tabular}{|c|c|c|c|c|c|}
	\hline
	$x$ & $-4$ & $-2$ & $0$ & $2$ & $4$ \\
	\hline
	$y = \dfrac{x^2}{2}$ & $8$ & $2$ & $0$ & $2$ & $8$ \\
	\hline
	\end{tabular}
	\end{center}
	Đồ thị $(P)$ là một đường parabol đi qua gốc toạ độ $O(0; 0)$, nhận trục tung $Oy$ làm trục đối xứng và nằm phía trên trục hoành.
	\item Điểm $M(x_M; y_M)$ thuộc $(P)$ có tung độ và hoành độ bằng nhau thoả mãn:
	$y_M = x_M \Leftrightarrow \dfrac{x_M^2}{2} = x_M \Leftrightarrow x_M^2 - 2x_M = 0 \Leftrightarrow x_M(x_M - 2) = 0 \Leftrightarrow \left[\begin{aligned} x_M &= 0 \Rightarrow y_M = 0 \\ x_M &= 2 \Rightarrow y_M = 2. \end{aligned}\right.$\\
	Vậy có hai điểm thoả mãn yêu cầu đề bài là $M_1(0; 0)$ và $M_2(2; 2)$.
\end{enumerate}
}
\end{bt}
